\section{Independence of torus closures}%
\label{sec:peps_torus_closure_independence}

On the torus a $G$-injective tensor is closed by group elements inserted along
the two seams (Definition~\ref{def:peps_torus_g_closure}). The closure vector
depends only on the pair-conjugacy class of the two elements, and this section
computes the overlaps of closure vectors for regular $G$-isometric tensors.

\begin{definition}[Pair-conjugacy classes of closures]%
    \label{def:peps_pair_conjugacy}
    \lean{TNLean.PEPS.PairConjugacyClass, TNLean.PEPS.pairConjugacyClass,
        TNLean.PEPS.pairConjugacyClass_eq_iff,
        TNLean.PEPS.commute_iff_of_pairConjugacyClass_eq,
        TNLean.PEPS.PairConjugacyClass.IsCommuting,
        TNLean.PEPS.PairConjugacyClass.isCommuting_pairConjugacyClass,
        TNLean.PEPS.CommutingPairConjugacyClass}
    \leanok
    Two pairs are equivalent if
    $(g,h)=(xg'x^{-1},xh'x^{-1})$ for some $x\in G$. Their equivalence
    classes are the pair-conjugacy classes of
    \cite[Definition~5.8]{Schuch2010PEPS}. Commutation is constant on each
    class, so the classes of commuting pairs are well-defined.
\end{definition}

\begin{lemma}[The torus vector depends only on the pair-conjugacy class]%
    \label{thm:peps_torus_pair_conjugacy}
    \lean{TNLean.PEPS.torusBondNetwork_conjugate, TNLean.PEPS.torusGClosure_conjugate,
        TNLean.PEPS.torusGClosure_eq_of_pairConjugacyClass_eq,
        TNLean.PEPS.torusGClosureClass, TNLean.PEPS.torusGClosureClass_pairConjugacyClass}
    \leanok
    \uses{def:peps_torus_g_closure, def:peps_pair_conjugacy}
    Suppose $A$ is invariant under the virtual group action. Then
    $\ket{\mathcal M(A\mid xgx^{-1},xhx^{-1})}
        =\ket{\mathcal M(A\mid g,h)}$ for every $x\in G$.
    Consequently $\ket{\mathcal M(A\mid C)}$ is well-defined for each
    pair-conjugacy class $C$ \cite[Definition~5.8]{Schuch2010PEPS}.
\end{lemma}

\begin{proof}\leanok
    \uses{lem:peps_torus_bond_network_gauge}
    In Lemma~\ref{lem:peps_torus_bond_network_gauge} take $P_v=U_x$ and
    $Q_v=U_{x^{-1}}$ at every site. The absorbed site tensor is unchanged
    by invariance. Each closure matrix $U_g$ or $U_h$ becomes
    $U_{xgx^{-1}}$ or $U_{xhx^{-1}}$, and each identity bond remains the
    identity. This proves the assertion for every coefficient in
    (\ref{eq:peps_torus_closure_coeff}).
\end{proof}

\begin{lemma}[Displacement of commuting closure seams]%
    \label{thm:peps_torus_commuting_closure_seams}
    \lean{TNLean.PEPS.torusHorizontalClosureAt_zero,
        TNLean.PEPS.torusHorizontalClosureAt,
        TNLean.PEPS.torusVerticalClosureAt,
        TNLean.PEPS.torusVerticalClosureAt_zero,
        TNLean.PEPS.torusBondNetwork_closureAt_eq_torusGClosure}
    \leanok
    \uses{def:peps_torus_g_closure}
    Let $A$ be invariant under the virtual group action and let $gh=hg$.
    For any column $c$ and row $r$, put
    \begin{align}
        H^c_v=\begin{cases}U_h,&v_x+1=c,\\\mathbb 1,&\text{otherwise},\end{cases}
        \qquad
        V^r_v=\begin{cases}U_g,&v_y+1=r,\\\mathbb 1,&\text{otherwise}.\end{cases}
    \end{align}
    Contracting $A$ with these bond operators gives
    $\ket{\mathcal M(A\mid g,h)}$, independently of $c$ and $r$.
    This is the displacement of commuting loops described after
    the string-deformation display of \cite{Schuch2010PEPS}, local source
    lines 1638--1644; neither isometry
    nor semi-regularity is needed. The identity holds on every nonempty
    rectangular torus, including dimensions equal to one.
\end{lemma}

\begin{proof}\leanok
    \uses{lem:peps_torus_bond_network_gauge}
    Choose coordinate representatives between zero and the corresponding
    circumference minus one. To move the horizontal seam, use the vertex
    change of basis $U_{h^{-1}}$ on columns preceding $c$ and the identity
    elsewhere. The adjacent changes cancel on all ordinary bonds, remove
    the old seam at zero, and insert $U_h$ at the new seam. They commute
    with every $U_g$ on the transverse seam. Similarly, use $U_g$ on rows
    preceding $r$ to move the vertical seam, with the inverse convention
    required by its orientation. Local invariance removes every absorbed
    change of basis. The network gauge identity therefore preserves each
    physical coefficient.
\end{proof}

\begin{definition}[Regular coordinates of a torus site]%
    \label{def:peps_torus_regular_coordinates}
    \lean{TNLean.PEPS.leftRegularMatrix,
        TNLean.PEPS.leftRegularMatrix_apply,
        TNLean.PEPS.toMatrix_leftRegular_eq_leftRegularMatrix,
        TNLean.PEPS.leftRegularMatrixEquiv,
        TNLean.PEPS.isSemiRegular_leftRegularMatrix,
        TNLean.PEPS.torusLegMatrix_leftRegularMatrix_apply,
        TNLean.PEPS.torusLegRep_leftRegularMatrix_apply,
        TNLean.PEPS.finFourArrowEquiv_smul,
        TNLean.PEPS.torusLegRep_leftRegularMatrix_conj,
        TNLean.PEPS.torusLegRep_leftRegularMatrix_unitary,
        TNLean.PEPS.toMatrix_averageMap_torusLegRep_leftRegularMatrix,
        TNLean.PEPS.toMatrix_siteMap_regular,
        TNLean.PEPS.siteMap_fourCoordinates}
    \leanok
    \uses{def:peps_torus_bond_network, def:peps_semiregular,
        thm:peps_semiregular_equivalence}
    In the group basis, $L_x\ket{k}=\ket{xk}$. The coordinate map
    $\sum_k c_k[k]\mapsto(c_k)_k$ intertwines the group-algebra regular
    representation with these permutation matrices. They are therefore
    semi-regular by Lemma~\ref{thm:peps_semiregular_equivalence}.
    The oriented four-leg
    action sends $(t,r,b,l)$ to $(xt,xr,xb,xl)$; its invariant projector
    has entries
    \begin{align}
        \Pi_{\eta,\theta}
        =\frac{1}{|G|}\sum_{x\in G}\mathbf{1}_{\eta=x\theta}.
    \end{align}
    Thus the four regular coordinates agree with the oriented virtual
    action in Definition~6.1 of \cite{Schuch2010PEPS}.
\end{definition}

\begin{lemma}[The local regular Gram identity]%
    \label{thm:peps_torus_regular_site_gram}
    \lean{TNLean.PEPS.IsGIsometric.torusRegularSite_translation,
        TNLean.PEPS.IsGIsometric.exists_torusRegularSiteGram}
    \leanok
    \uses{def:peps_torus_regular_coordinates, def:peps_g_injective}
    A regular $G$-isometric site tensor $A$ is invariant under common
    translation of its four virtual labels. There is $c>0$ such that
    \begin{align}
        \sum_s\overline{A_s(\eta)}A_s(\theta)
        =\frac{c}{|G|}\sum_{x\in G}\mathbf{1}_{\eta=x\theta}.
    \end{align}
    The factor $c$ records the normalization of the site tensor.
\end{lemma}

\begin{proof}\leanok
    Expressing $A^\dagger A=c\Pi$ in the orthonormal group basis and substituting
    the matrix entries of the invariant projector yields the formula.
\end{proof}

\begin{lemma}[Access to the virtual invariant space]%
    \label{thm:peps_torus_accessible_virtual_site}
    \lean{TNLean.PEPS.IsGIsometric.exists_accessibleVirtualSite}
    \leanok
    If $A^\dagger A=c\Pi$ with $c>0$, then $J=c^{-1/2}A^\dagger$
    satisfies $JA=\sqrt{c}\Pi$ and preserves inner products on
    $\operatorname{ran}A$. It identifies the physical support with
    the invariant virtual space, as in Observation~6.4 of
    \cite{Schuch2010PEPS}.
\end{lemma}

\begin{proof}\leanok
    \uses{thm:peps_torus_regular_site_gram}
    For $u,v$ in the virtual space, $JA u=\sqrt{c}\Pi u$ and
    $\langle Au,Av\rangle=c\langle\Pi u,\Pi v\rangle$.
\end{proof}

\begin{lemma}[Expansion of a torus overlap]%
    \label{thm:peps_torus_overlap_expansion}
    \lean{TNLean.PEPS.sum_star_weightedSiteProduct_mul,
        TNLean.PEPS.torusBondSiteLabels,
        TNLean.PEPS.torusBondWeight,
        TNLean.PEPS.sum_star_torusSiteProduct_mul,
        TNLean.PEPS.sum_star_torusBondNetwork_mul,
        TNLean.PEPS.permMatrixHom_apply_eq_ite,
        TNLean.PEPS.torusPermutationSiteLabels,
        TNLean.PEPS.torusBondNetwork_perm,
        TNLean.PEPS.sum_star_torusBondNetwork_perm_mul}
    \leanok
    \uses{def:peps_torus_bond_network}
    Let $W(\beta)$ denote the product of the bond matrix entries and
    let $\eta_v(\beta)$ be the four labels read by site $v$.
    The overlap of two networks with site tensors $A_v,B_v$ is
    \begin{align}
        \sum_{\beta,\gamma}\overline{W_A(\beta)}W_B(\gamma)
        \prod_v\left(\sum_s
        \overline{A_{v,s}(\eta_v(\beta))}
        B_{v,s}(\eta_v(\gamma))\right).
    \end{align}
    For permutation bonds, the head label is the permutation of the
    tail label. Hence one free label per bond suffices in each network.
\end{lemma}

\begin{proof}\leanok
    Expanding both contractions, interchanging the finite sums, and summing each
    physical index independently yields the product formula, while a permutation matrix has entry one
    exactly when its row is the image of its column.
\end{proof}

\begin{lemma}[Counting compatible regular bonds]%
    \label{thm:peps_regular_bond_intertwining_count}
    \lean{TNLean.PEPS.regularBond_labels_iff_intertwining,
        TNLean.PEPS.sum_regularBond_labels_eq_card,
        TNLean.PEPS.sum_two_regularBond_labels_eq_card,
        TNLean.PEPS.torusRegularSiteLabels,
        TNLean.PEPS.torusRegularSiteLabels_translation_iff,
        TNLean.PEPS.sum_torusRegularSiteLabels_translation_eq_card,
        TNLean.PEPS.sum_torusClosureSiteLabels_translation_eq_card}
    \leanok
    \uses{thm:peps_torus_overlap_expansion}
    On a bond $e$ oriented from $s(e)$ to $t(e)$, let $u_e,w_e$ be
    the group elements in the two networks. Vertex translations $q_v$
    are compatible with this bond precisely when
    \begin{align}
        u_e q_{s(e)}=q_{t(e)}w_e.
    \end{align}
    If the equations hold, each bond has one independent group label;
    otherwise there are no compatible labels. For the two bond families
    on a torus with $N$ sites, the count is $|G|^{2N}$.
\end{lemma}

\begin{proof}\leanok
    The tail labels satisfy $\eta_e=q_{s(e)}\theta_e$. Substituting this
    into $u_e\eta_e=q_{t(e)}w_e\theta_e$ and cancelling $\theta_e$
    gives the intertwining equation. Every $\theta_e$ is then free.
\end{proof}

\begin{lemma}[Translation compatibility across the torus seams]%
    \label{thm:peps_torus_closure_compatibility}
    \lean{TNLean.PEPS.torusHorizontalClosure_eq_map_element,
        TNLean.PEPS.torusHorizontalClosureElement,
        TNLean.PEPS.torusVerticalClosureElement,
        TNLean.PEPS.torusVerticalClosure_eq_map_element,
        TNLean.PEPS.IsTorusClosureCompatible,
        TNLean.PEPS.IsTorusClosureCompatible.eq_origin,
        TNLean.PEPS.isTorusClosureCompatible_iff_exists_intertwiner,
        TNLean.PEPS.sum_torusClosureCompatible_eq_sum_intertwiner}
    \leanok
    \uses{def:peps_torus_g_closure,
        thm:peps_regular_bond_intertwining_count}
    For closures $(g,h)$ and $(g',h')$, a compatible vertex translation
    is constant, $q_v=x$, and its value satisfies
    \begin{align}
        hx=xh',\qquad gx=xg'.
        \label{eq:peps_closure_seam_intertwining}
    \end{align}
    Summing over compatible vertex translations therefore gives exactly
    the number of simultaneous intertwiners. This holds also when a
    torus circumference is one.
\end{lemma}

\begin{proof}\leanok
    The nonseam horizontal and vertical bonds identify all translations
    in a rectangular fundamental domain. The two seams give the equations
    (\ref{eq:peps_closure_seam_intertwining}). For a circumference of one the corresponding seam is
    already a loop at each site, imposing the same equation.
\end{proof}

\begin{lemma}[Overlap of regular isometric torus closures]%
    \label{thm:peps_regular_torus_closure_overlap}
    \lean{TNLean.PEPS.prod_torusSiteGram_eq_sum_translation,
        TNLean.PEPS.sum_star_torusBondNetwork_regularClosure_mul,
        TNLean.PEPS.sum_star_torusGClosure_regular_mul,
        TNLean.PEPS.IsGIsometric.exists_torusGClosure_overlap}
    \leanok
    \uses{thm:peps_torus_overlap_expansion, thm:peps_torus_closure_compatibility}
    Let $A_v^\dagger A_v=c_v\Pi_v$ be the local regular isometry identities
    on a torus with $N$ sites, with $c_v>0$. For the closure vectors,
    \begin{align}
        \langle\psi_{g,h},\psi_{g',h'}\rangle
        =\left(\prod_v c_v\right)|G|^N
        \#\{x\in G:hx=xh',\ gx=xg'\}.
    \end{align}
    For a uniform tensor the prefactor is $c^N|G|^N$. Neither closure
    pair is required to commute.
\end{lemma}

\begin{proof}\leanok
    \uses{thm:peps_torus_regular_site_gram, thm:peps_finite_indicator_sum}
    Expanding the overlap into local Gram kernels expresses their product as the
    sum over vertex translations with prefactor
    $(\prod_v c_v)|G|^{-N}$. For a compatible translation, the two bond
    families contribute $|G|^{2N}$, leaving the stated prefactor. By
    Lemma~\ref{thm:peps_torus_closure_compatibility}, the remaining sum
    counts the simultaneous intertwiners.
\end{proof}

\begin{lemma}[Independent simultaneous-conjugacy operator sums]%
    \label{thm:peps_pair_conjugacy_operator_independence}
    \lean{TNLean.PEPS.pairConjugacyOperator_eq_of_pairConjugacyClass_eq,
        TNLean.PEPS.pairConjugacyOperator,
        TNLean.PEPS.pairConjugacyClassOperator,
        TNLean.PEPS.pairConjugacyClassOperator_pairConjugacyClass,
        TNLean.PEPS.pairOperatorDeltaPairing,
        TNLean.PEPS.pairOperatorDeltaPairing_tmul,
        TNLean.PEPS.pairOperatorDeltaPairing_pairConjugacyOperator,
        TNLean.PEPS.pairOperatorDeltaPairing_eq_zero_of_pairConjugacyClass_ne,
        TNLean.PEPS.pairOperatorDeltaPairing_self_eq_card_centralizer,
        TNLean.PEPS.pairOperatorDeltaPairing_self_ne_zero,
        TNLean.PEPS.linearIndependent_pairConjugacyClassOperator,
        TNLean.PEPS.pairConjugacyClassKroneckerOperator,
        TNLean.PEPS.pairConjugacyClassKroneckerOperator_pairConjugacyClass,
        TNLean.PEPS.linearIndependent_pairConjugacyClassKroneckerOperator}
    \leanok
    \uses{def:peps_pair_conjugacy, def:peps_semiregular, def:peps_delta_operator}
    Let $U$ be a semi-regular representation of a finite group $G$ on a
    finite-dimensional complex vector space $E$. Define
    \begin{align}
        \Omega_{g,h}
         & =\sum_{x\in G}U_{xhx^{-1}}\otimes U_{xgx^{-1}},
        \label{eq:peps_pair_operator_sum}                  \\
        \Phi_{g,h}(A\otimes B)
         & =\tr[U_{h^{-1}}A\Delta]\,
        \tr[U_{g^{-1}}B\Delta].
        \label{eq:peps_pair_operator_dual}
    \end{align}
    The operator $\Omega_{g,h}$ depends only on the simultaneous-conjugacy
    class $C$ of $(g,h)$; denote it by $\Omega_C$. The family
    $(\Omega_C)_C$ is linearly independent. The detecting functionals satisfy
    \begin{align}
        \Phi_{g,h}(\Omega_{g',h'})
         & =\bigl|\{x\in G:
        g=xg'x^{-1},\ h=xh'x^{-1}\}\bigr|,
        \label{eq:peps_pair_operator_count} \\
        \Phi_{g,h}(\Omega_{g,h})
         & =|Z_G(g)\cap Z_G(h)|>0.
        \label{eq:peps_pair_operator_diagonal}
    \end{align}
    In a coordinate basis, the corresponding matrices
    $\sum_x U_{xhx^{-1}}\otimes_{\mathrm K}U_{xgx^{-1}}$ are likewise
    well-defined on classes and linearly independent. Here
    $\otimes_{\mathrm K}$ denotes the Kronecker product.
    Restriction to commuting pair classes gives the operator-independence
    argument in \cite[Theorem~5.9]{Schuch2010PEPS}.
\end{lemma}

\begin{proof}\leanok
    \uses{def:peps_pair_conjugacy, thm:peps_semiregular_trace}
    Replacing $(g,h)$ by $(ygy^{-1},yhy^{-1})$ in
    (\ref{eq:peps_pair_operator_sum}) and substituting $x\mapsto xy$
    leaves the sum unchanged. By Lemma~\ref{thm:peps_semiregular_trace},
    each factor of (\ref{eq:peps_pair_operator_dual}) evaluates to a
    Kronecker delta on a representation operator, giving
    (\ref{eq:peps_pair_operator_count}). Its value is zero on distinct
    classes, while (\ref{eq:peps_pair_operator_diagonal}) is nonzero
    because the common centralizer contains the identity. Applying
    $\Phi_{g,h}$ to a finite relation $\sum_C\lambda_C\Omega_C=0$
    therefore gives $\lambda_C|Z_G(g)\cap Z_G(h)|=0$ for the class of
    $(g,h)$, and hence every coefficient vanishes. Finally, the canonical
    isomorphism
    $\operatorname{End}(E)\otimes\operatorname{End}(E)
        \cong\operatorname{End}(E\otimes E)$, followed by passage to the
    tensor-product coordinate basis, takes simple tensors to Kronecker
    matrices and preserves linear independence.
\end{proof}
